\chapter{Number-Theoretic Algorithms}

\SetKw{KwDownTo}{downto}

In Chapter~29, we manipulated polynomials and signals using complex roots of unity. Here we turn from continuous rotations to discrete arithmetic on whole numbers: which integers divide which, what is left over when they do not, and how arithmetic behaves when we retain only the remainder. This is \emph{number theory}, the mathematical foundation of modern cryptography.

For most of this book, data items have been treated as opaque tokens. In this chapter we open integers up to inspect their internal structure. We begin with Euclid's greatest common divisor algorithm, extend it to solve modular linear equations and compute multiplicative inverses, use repeated squaring for fast modular exponentiation, and combine these tools into the Miller--Rabin randomized primality test and the RSA cryptosystem.

Throughout, we track both arithmetic operations and bit-level costs for $\beta$-bit operands, since what makes these algorithms useful is that they stay efficient even when numbers span hundreds of digits.

\chapterline

\section{Divisibility, GCD, and the Idea of a Remainder}
\label{ch30:divisibility}

Let us begin where every number-theory course begins, with the relation of one
integer \emph{dividing} another.

\begin{definition}[Divisibility]
\label{ch30:def-divides}
For integers $a$ and $d$ with $d \neq 0$, we say $d$ \emph{divides} $a$, written
$d \mid a$, if there is an integer $k$ with $a = k\,d$. We call $d$ a
\emph{divisor} of $a$ and $a$ a \emph{multiple} of $d$. If $d$ does not divide
$a$ we write $d \nmid a$.
\end{definition}

So $3 \mid 12$ because $12 = 4 \cdot 3$, but $3 \nmid 13$. Every integer $d$
divides $0$ (take $k = 0$), and $1$ divides everything.

Everything that follows depends on one more fact: when $d$ does \emph{not}
evenly divide $a$, there is a unique leftover.

\begin{theorem}[Division theorem]
\label{ch30:thm-division}
For any integer $a$ and any positive integer $n$, there exist unique integers
$q$ (the \emph{quotient}) and $r$ (the \emph{remainder}) such that
\[
  a = q\,n + r \qquad\text{and}\qquad 0 \le r < n.
\]
We write $q = \floor{a/n}$ and $r = a \bmod n$.
\end{theorem}

You have known this since grade school as ``long division with a remainder.''
The remainder $a \bmod n$ is always dragged into the range $\{0, 1, \dots,
n-1\}$, even when $a$ is negative: $(-7) \bmod 3 = 2$, because $-7 = (-3)\cdot 3 +
2$. Hold onto that; sign conventions for $\bmod$ trip up almost everyone at
least once.

\subsection{The greatest common divisor}

\begin{definition}[GCD]
\label{ch30:def-gcd}
The \emph{greatest common divisor} of two integers $a$ and $b$, not both zero,
written $\gcd(a,b)$, is the largest integer that divides both. We set
$\gcd(0,0) = 0$ by convention. If $\gcd(a,b) = 1$ we say $a$ and $b$ are
\emph{relatively prime} (or \emph{coprime}).
\end{definition}

How would you compute $\gcd(48, 36)$ if you had never heard of Euclid? You might
factor: $48 = 2^4 \cdot 3$ and $36 = 2^2 \cdot 3^2$, take the smaller power of
each shared prime, $2^2 \cdot 3 = 12$. That is correct, but it
does not scale: \emph{factoring large numbers is believed to be hard}.
Nobody knows a fast way to factor a 600-digit number, and the security of RSA
leans on exactly that difficulty. So if computing $\gcd$ required factoring, we
would be stuck. The surprise is that you can compute $\gcd$
without factoring at all, and very fast, using an idea about 2300 years old.

\begin{intuitionbox}
Picture a floor that is $48$ by $36$ inches, and you want to tile it with the
largest possible square tiles, all the same size, no cutting. The side length of
that biggest tile is exactly $\gcd(48,36)$. Here is Euclid's insight as a
physical act. Lay down as many $36 \times 36$ squares along the long side as you
can -- you fit one, with a $12 \times 36$ strip left over. Now the leftover strip
is itself a smaller tiling problem: tile a $36 \times 12$ region. The key fact is
that the biggest tile for the original floor is the same as the biggest tile for
the leftover strip. The problem got smaller but the answer did not change. Keep
peeling off the leftover and you eventually reach a strip that tiles perfectly --
and that final side length is your answer.
\end{intuitionbox}

Figure~\ref{ch30:fig-tiling} draws the first peel of that process for a
$48 \times 36$ floor: one $36 \times 36$ square fits along the long side, leaving
a $12 \times 36$ strip that becomes the next, smaller tiling problem.

\begin{figure}[H]
\centering
\begin{tikzpicture}[scale=0.072]
  % full floor 48 wide x 36 tall
  \draw[thick] (0,0) rectangle (48,36);
  % the 36x36 square peeled off
  \fill[sourceblue!18] (0,0) rectangle (36,36);
  \draw[sourceblue, thick] (0,0) rectangle (36,36);
  % the leftover 12x36 strip
  \fill[quotientorange!22] (36,0) rectangle (48,36);
  \draw[quotientorange, thick] (36,0) rectangle (48,36);
  % dimension labels
  \node at (18,18) {\footnotesize $36\times 36$ square};
  \node[rotate=90] at (42,18) {\scriptsize $12\times 36$ left};
  \draw[<->, >=Stealth] (0,-3) -- (48,-3) node[midway,below,font=\scriptsize]{$48$};
  \draw[<->, >=Stealth] (-3,0) -- (-3,36) node[midway,left,font=\scriptsize]{$36$};
\end{tikzpicture}
\caption{One step of Euclid as tiling: $\gcd(48,36)=\gcd(36,12)$. The blue
$36\times 36$ square is the largest that fits along the long side; the orange
$12\times 36$ leftover is the next subproblem, with the same answer.}
\label{ch30:fig-tiling}
\end{figure}

The picture rests on one short lemma.

\begin{theorem}[Euclid's reduction]
\label{ch30:thm-euclid}
For any integer $a$ and positive integer $b$,
\[
  \gcd(a,b) = \gcd(b,\; a \bmod b).
\]
\end{theorem}

\begin{formalbox}
\textit{Proof.} Write $a = q b + r$ with $r = a \bmod b$. We show the two pairs
$(a,b)$ and $(b,r)$ have \emph{exactly the same set} of common divisors; then in
particular their greatest common divisors coincide.

Suppose $d \mid a$ and $d \mid b$. Then $d \mid (a - qb) = r$, so $d$ divides
both $b$ and $r$. Conversely, suppose $d \mid b$ and $d \mid r$. Then $d \mid
(qb + r) = a$, so $d$ divides both $a$ and $b$. The common-divisor sets are
equal, hence so are their maxima. \hfill$\square$
\end{formalbox}

Euclid's algorithm simply applies this fact over and over. We replace
the pair $(a,b)$ by the smaller pair $(b, a \bmod b)$ and repeat. The second
coordinate strictly shrinks each round and stays non-negative, so we must reach
$0$. When the second coordinate is $0$, the first coordinate is the answer,
because $\gcd(g, 0) = g$.

\begin{algorithm}[H]
\caption{\textsc{Euclid}$(a,b)$ -- greatest common divisor}
\label{ch30:alg-euclid}
\KwIn{Non-negative integers $a, b$, not both zero.}
\KwOut{$\gcd(a,b)$.}
\While{$b \neq 0$}{
  $r \leftarrow a \bmod b$\;
  $a \leftarrow b$\;
  $b \leftarrow r$\;
}
\Return $a$\;
\end{algorithm}

\begin{example}[Tracing $\gcd(252,198)$]
\label{ch30:ex-euclid}
Each line replaces $(a,b)$ with $(b,\,a \bmod b)$.
\begin{center}
\begin{tabular}{r r r r}
\toprule
$a$ & $b$ & $q = \floor{a/b}$ & $r = a \bmod b$ \\
\midrule
$252$ & $198$ & $1$ & $54$ \\
$198$ & $54$  & $3$ & $36$ \\
$54$  & $36$  & $1$ & $18$ \\
$36$  & $18$  & $2$ & $0$ \\
$18$  & $0$   & --- & --- \\
\bottomrule
\end{tabular}
\end{center}
The loop stops when $b = 0$, returning $a = 18$. And indeed $252 = 14 \cdot 18$
and $198 = 11 \cdot 18$, with $\gcd(14,11) = 1$, so $18$ is right.
\end{example}

\subsection{How fast is it? Enter Fibonacci}

Each iteration is cheap, so the only question is: how many iterations can there
be? The answer ties this ancient algorithm to a sequence you have met before.

\begin{theorem}[Lam\'e's theorem, informal]
\label{ch30:thm-lame}
If $a > b \ge 1$ and the call \textsc{Euclid}$(a,b)$ performs $k \ge 1$
recursive reductions, then $a \ge F_{k+2}$ and $b \ge F_{k+1}$, where $F_i$ is the
$i$-th Fibonacci number ($F_0 = 0$, $F_1 = 1$, $F_{i} = F_{i-1} + F_{i-2}$). Equivalently, the number of reductions is $O(\log b)$.
\end{theorem}

\begin{intuitionbox}
Why Fibonacci? Each reduction step replaces $(a,b)$ with $(b, a \bmod b)$. To
make the algorithm work as \emph{hard} as possible -- the most steps for the
smallest inputs -- you want the remainder $a \bmod b$ to be as large as it can be
while still being a genuine reduction, and you want the quotient to be the
smallest possible, namely $1$. A quotient of $1$ at every single step means
$a = b + r$, i.e.\ each number is the \emph{sum} of the next two. That recurrence,
``each term is the sum of the previous two,'' is the definition of the Fibonacci
numbers. So consecutive Fibonacci numbers are the worst case: $\gcd(F_{n+1},
F_n)$ takes the maximum number of steps.
\end{intuitionbox}

Because $F_k$ grows like $\phi^k$ with $\phi = (1+\sqrt5)/2 \approx 1.618$, the
condition $a \ge F_{k+2}$ forces $k = O(\log a)$. Concretely, the number of
divisions is at most about $\log_\phi a \approx 1.44 \log_2 a$ -- under five
divisions per decimal digit. That is why Euclid's algorithm is fast even on
huge inputs: a 600-digit number triggers only a few thousand divisions.

\begin{complexitybox}
\textbf{Euclid's algorithm.} On inputs $a > b \ge 0$:
\begin{itemize}
  \item \emph{Number of divisions:} $O(\log b)$, worst case at consecutive
        Fibonacci numbers.
  \item \emph{Arithmetic operations:} $O(\log b)$ modulo operations.
  \item \emph{Bit complexity:} each division of $\beta$-bit numbers costs
        $O(\beta^2)$ with schoolbook arithmetic, giving an easy $O(\beta^3)$
        bound; a sharper accounting (the quotients shrink the numbers) brings the
        true cost down to $O(\beta^2)$, and fast multiplication improves it
        further.
\end{itemize}
\end{complexitybox}

\begin{trapbox}
Do \emph{not} ``optimize'' Euclid by factoring first. Factoring is exponentially
slower than the GCD itself for large numbers. Euclid's method is fast precisely
because it avoids factoring entirely. Also: the algorithm needs $a \bmod b$,
not $a - b$ -- repeated subtraction is correct but can take $a/b$ steps in a
single ``level,'' which is catastrophic when, say, $a = 10^9$ and $b = 1$.
\end{trapbox}

\section{The Extended Euclidean Algorithm}
\label{ch30:extended}

Euclid's algorithm tells you \emph{what} $\gcd(a,b)$ is. With a little extra
bookkeeping it also expresses $\gcd(a,b)$ as a combination of $a$ and $b$, and
that extended version is the most useful tool in this chapter.

\begin{theorem}[B\'ezout's identity]
\label{ch30:thm-bezout}
For any integers $a$ and $b$, not both zero, there exist integers $x$ and $y$
with
\[
  a x + b y = \gcd(a,b).
\]
In fact, $\gcd(a,b)$ is the \emph{smallest} positive integer expressible as an
integer combination $ax + by$.
\end{theorem}

\begin{intuitionbox}
Geometrically, consider all integers you can reach on the number line by
taking integer combinations of step sizes $a$ and $b$:
\[
L = \{ax + by : x, y \in \Z\} .
\]
This set forms a discrete regular grid (a one-dimensional sub-lattice or ideal)
on the number line. Because adding or subtracting any two numbers in $L$ produces
another number in $L$, the entire lattice consists precisely of all integer
multiples of its smallest positive element $d$. Any common divisor of $a$ and
$b$ must divide every number in $L$, including $d$; conversely, since $a = a\cdot
1 + b\cdot 0 \in L$ and $b = a\cdot 0 + b\cdot 1 \in L$, that minimal positive
generator $d$ must divide both $a$ and $b$. Therefore, $d$ is exactly
$\gcd(a,b)$.

B\'ezout's identity says that the greatest common divisor is also an integer
linear combination of $a$ and $b$: you can always express $18$ as a
combination of $252$ and $198$, and indeed $252 \cdot 4 + 198 \cdot (-5) = 1008
- 990 = 18$. The integers $x$ and $y$ are the \emph{B\'ezout coefficients}, and
extracting them is the job of the extended Euclidean algorithm.
\end{intuitionbox}

Here is how the algorithm extracts them while running Euclid: write each
remainder as a combination of the original $a$ and $b$. At the start, $a =
1\cdot a + 0 \cdot b$ and $b = 0\cdot a + 1\cdot b$. Every later remainder is
built by subtracting a multiple of one previous line from another: $r = a - qb$.
Since both inputs to that subtraction are already combinations of the original
$a$ and $b$, so is the result. The combinations ride along for free on top of
the ordinary GCD computation. When you reach the GCD, you have its combination
too.

The cleanest way to compute the coefficients is the recursive form. Suppose the
recursive call on $(b, a \bmod b)$ has already handed back a triple $(d, x', y')$
with $b x' + (a \bmod b) y' = d$. We want $(d, x, y)$ with $a x + b y = d$.
Substitute $a \bmod b = a - \floor{a/b}\, b$:
\[
  d = b x' + (a - \floor{a/b}\,b)\,y'
    = a\,y' + b\bigl(x' - \floor{a/b}\,y'\bigr).
\]
Matching coefficients, $x = y'$ and $y = x' - \floor{a/b}\,y'$.

\begin{algorithm}[H]
\caption{\textsc{Extended-Euclid}$(a,b)$}
\label{ch30:alg-exteuclid}
\KwIn{Non-negative integers $a,b$.}
\KwOut{Triple $(d,x,y)$ with $d=\gcd(a,b)$ and $ax+by=d$.}
\If{$b = 0$}{
  \Return $(a,\,1,\,0)$\;
}
$(d,\, x',\, y') \leftarrow \textsc{Extended-Euclid}(b,\; a \bmod b)$\;
$(d,\, x,\, y) \leftarrow (d,\; y',\; x' - \floor{a/b}\cdot y')$\;
\Return $(d,x,y)$\;
\end{algorithm}

\begin{example}[Coefficients for $\gcd(252,198)$]
\label{ch30:ex-exteuclid}
We recurse down to the base case, then unwind, building the coefficients on the
way back up. The downward pass is exactly Example~\ref{ch30:ex-euclid}; here is
the upward pass, each line using $x = y'$ and $y = x' - \floor{a/b}y'$.
\begin{center}
\begin{tabular}{r r r r r r}
\toprule
$a$ & $b$ & $\floor{a/b}$ & $d$ & $x$ & $y$ \\
\midrule
$18$  & $0$   & ---  & $18$ & $1$  & $0$  \\
$36$  & $18$  & $2$  & $18$ & $0$  & $1$  \\
$54$  & $36$  & $1$  & $18$ & $1$  & $-1$ \\
$198$ & $54$  & $3$  & $18$ & $-1$ & $4$  \\
$252$ & $198$ & $1$  & $18$ & $4$  & $-5$ \\
\bottomrule
\end{tabular}
\end{center}
The top line of the original call returns $(18, 4, -5)$, and we verify $252\cdot
4 + 198\cdot(-5) = 1008 - 990 = 18$. The B\'ezout coefficients took the
same number of steps as the plain GCD.
\end{example}

Figure~\ref{ch30:fig-exteuclid} illustrates the two-way dataflow of the
extended Euclidean algorithm: the downward pass discovers the quotients and
remainders via ordinary division, while the upward pass reassembles the B\'ezout
coefficients from the bottom up.

\begin{figure}[H]
\centering
\begin{tikzpicture}[
  xscale=1.0, yscale=1.0, >=Stealth,
  callbox/.style={draw=sourceblue!70, rounded corners=3pt, fill=sourceblue!8, text width=4.4cm, inner sep=4pt, font=\scriptsize},
  retbox/.style={draw=targetgreen!70, rounded corners=3pt, fill=targetgreen!8, text width=4.8cm, inner sep=4pt, font=\scriptsize},
  desc/.style={->, line width=1.1pt, sourceblue},
  asc/.style={->, line width=1.1pt, targetgreen!80!black}]

  % Lane headers
  \node[font=\footnotesize\bfseries, sourceblue!90!black] at (2.2, 0.7) {(a) Downward Division Steps};
  \node[font=\footnotesize\bfseries, targetgreen!90!black] at (7.6, 0.7) {(b) Upward B\'ezout Assembly};

  % Downward calls (left lane)
  \node[callbox] (c0) at (2.2, 0) {
    \textbf{Call 0:} \textsc{Ext-Euclid}$(252, 198)$\\
    $252 = 1 \cdot 198 + 54$ \hfill $(q=1, r=54)$
  };

  \node[callbox] (c1) at (2.2, -1.25) {
    \textbf{Call 1:} \textsc{Ext-Euclid}$(198, 54)$\\
    $198 = 3 \cdot 54 + 36$ \hfill $(q=3, r=36)$
  };

  \node[callbox] (c2) at (2.2, -2.5) {
    \textbf{Call 2:} \textsc{Ext-Euclid}$(54, 36)$\\
    $54 = 1 \cdot 36 + 18$ \hfill $(q=1, r=18)$
  };

  \node[callbox] (c3) at (2.2, -3.75) {
    \textbf{Call 3:} \textsc{Ext-Euclid}$(36, 18)$\\
    $36 = 2 \cdot 18 + 0$ \hfill $(q=2, r=0)$
  };

  \node[callbox, draw=trapred!60, fill=trapred!8] (c4) at (2.2, -5.0) {
    \textbf{Call 4 (Base):} \textsc{Ext-Euclid}$(18, 0)$\\
    $b=0 \implies d=18$, return $(18, 1, 0)$
  };

  % Downward recursion arrows
  \draw[desc] (c0.south west) ++(0.3,0) -- node[left, font=\tiny]{$q=1$} ++(0,-0.45);
  \draw[desc] (c1.south west) ++(0.3,0) -- node[left, font=\tiny]{$q=3$} ++(0,-0.45);
  \draw[desc] (c2.south west) ++(0.3,0) -- node[left, font=\tiny]{$q=1$} ++(0,-0.45);
  \draw[desc] (c3.south west) ++(0.3,0) -- node[left, font=\tiny]{$q=2$} ++(0,-0.45);

  % Upward returns (right lane)
  \node[retbox, draw=trapred!60, fill=trapred!8] (r4) at (7.6, -5.0) {
    \textbf{Return Base:} $(d,x,y) = (18, 1, 0)$\\
    Check: $18(1) + 0(0) = 18$
  };

  \node[retbox] (r3) at (7.6, -3.75) {
    $x=0,\; y=1-2(0)=1 \implies (18, 0, 1)$\\
    Check: $36(0) + 18(1) = 18$
  };

  \node[retbox] (r2) at (7.6, -2.5) {
    $x=1,\; y=0-1(1)=-1 \implies (18, 1, -1)$\\
    Check: $54(1) + 36(-1) = 18$
  };

  \node[retbox] (r1) at (7.6, -1.25) {
    $x=-1,\; y=1-3(-1)=4 \implies (18, -1, 4)$\\
    Check: $198(-1) + 54(4) = 18$
  };

  \node[retbox, line width=0.9pt, draw=targetgreen!80!black] (r0) at (7.6, 0) {
    \textbf{Final Result:} $(d,x,y) = (18, 4, -5)$\\
    Check: $252(4) + 198(-5) = 18$
  };

  % Base horizontal transition
  \draw[->, line width=1.1pt, trapred] (c4.east) -- node[above, font=\tiny\bfseries, trapred]{init} (r4.west);

  % Upward unwinding arrows
  \draw[asc] (r4.north east) ++(-0.4,0) -- node[right, font=\tiny]{unwind} ++(0,0.45);
  \draw[asc] (r3.north east) ++(-0.4,0) -- node[right, font=\tiny]{unwind} ++(0,0.45);
  \draw[asc] (r2.north east) ++(-0.4,0) -- node[right, font=\tiny]{unwind} ++(0,0.45);
  \draw[asc] (r1.north east) ++(-0.4,0) -- node[right, font=\tiny]{unwind} ++(0,0.45);

\end{tikzpicture}
\caption{The Extended Euclidean algorithm on $(252,198)$. (a) The downward pass
executes standard Euclidean division $a = qb + r$, strictly reducing inputs down
to the base case $(18,0)$. (b) The upward pass unwinds via $x = y'$ and
$y = x' - qy'$, producing the B\'ezout linear combination $252(4) + 198(-5) = 18$
at the root with no extra division steps.}
\label{ch30:fig-exteuclid}
\end{figure}

\begin{focusbox}
The extended algorithm gives you, for free, an integer $x$ such that $ax \equiv
\gcd(a,b) \pmod{b}$. When $\gcd(a,b)=1$, that $x$ is the \emph{multiplicative inverse} of
$a$ modulo $b$. Nearly every major tool in the rest of this chapter -- solving
modular linear equations, the Chinese Remainder Theorem, and RSA decryption --
builds directly on computing modular inverses. If you remember one algorithm from this chapter, make
it \textsc{Extended-Euclid}.
\end{focusbox}

\section{Modular Arithmetic: Computing on a Clock}
\label{ch30:modular}

A clock face is a familiar example. Seven hours after $9$ o'clock is
not $16$ o'clock; it is $4$, because we wrap around at $12$. We threw away the
``how many times around'' information and kept only the position on the dial.
Modular arithmetic does exactly this, with $n$ positions instead of $12$.
Figure~\ref{ch30:fig-clock} shows the wrap-around for $9 + 7 \equiv 4 \pmod{12}$.

\begin{figure}[H]
\centering
\begin{tikzpicture}[scale=0.95, >=Stealth]
  % Panel (a): Addition on Z_12
  \begin{scope}[shift={(0,0)}]
    \def\R{1.45}
    \node[anchor=south, font=\small\bfseries, black!85] at (0, 2.2) {(a) Addition on $\Z_{12}$: $9 + 7 \equiv 4$};
    \draw[thick, black!50] (0,0) circle (\R);
    \foreach \k in {0,...,11}{
      \pgfmathsetmacro\ang{90 - \k*30}
      \node[font=\tiny] at (\ang:\R+0.25) {\k};
      \fill (\ang:\R) circle (0.6pt);
    }
    % start at 9
    \pgfmathsetmacro\angA{90 - 9*30}
    \draw[sourceblue, very thick, -{Stealth}] (0,0) -- (\angA:\R-0.12);
    \node[sourceblue, font=\tiny\bfseries] at (-0.65, 0.4) {start $9$};

    % end at 4
    \pgfmathsetmacro\angB{90 - 4*30}
    \draw[targetgreen, very thick, -{Stealth}] (0,0) -- (\angB:\R-0.12);
    \node[targetgreen!80!black, font=\tiny\bfseries] at (0.55, -0.6) {end $4$};

    % arc showing +7 clockwise from 9 to 4
    \draw[quotientorange, thick, -{Stealth}] (\angA:\R+0.5)
      arc[start angle=\angA, end angle=\angB-360, radius=\R+0.5];
    \node[quotientorange!90!black, font=\footnotesize\bfseries] at (-1.85, 1.05) {$+7$};
  \end{scope}

  % Panel (b): Multiplication & Inverses on Z_7
  \begin{scope}[shift={(6.2,0)}]
    \def\R{1.45}
    \node[anchor=south, font=\small\bfseries, black!85] at (0, 2.2) {(b) Multiplication \& Inverses on $\Z_7$};
    \draw[thick, black!50] (0,0) circle (\R);
    \foreach \k in {0,...,6}{
      \pgfmathsetmacro\ang{90 - \k*(360/7)}
      \node[font=\tiny] at (\ang:\R+0.25) {\k};
      \fill (\ang:\R) circle (0.6pt);
    }

    % Orbit of multiplying by 3: 1 -> 3 -> 2 -> 6 -> 4 -> 5 -> 1
    \pgfmathsetmacro\aOne{90 - 1*(360/7)}
    \pgfmathsetmacro\aThree{90 - 3*(360/7)}
    \pgfmathsetmacro\aTwo{90 - 2*(360/7)}
    \pgfmathsetmacro\aSix{90 - 6*(360/7)}
    \pgfmathsetmacro\aFour{90 - 4*(360/7)}
    \pgfmathsetmacro\aFive{90 - 5*(360/7)}

    % Multiplication paths inside the dial
    \draw[->, black!35, dashed] (\aOne:\R-0.2) -- (\aThree:\R-0.2);
    \draw[->, black!35, dashed] (\aThree:\R-0.2) -- (\aTwo:\R-0.2);
    \draw[->, black!35, dashed] (\aTwo:\R-0.2) -- (\aSix:\R-0.2);
    \draw[->, black!35, dashed] (\aSix:\R-0.2) -- (\aFour:\R-0.2);
    \draw[->, black!35, dashed] (\aFour:\R-0.2) -- (\aFive:\R-0.2);
    
    % Key inverse arrow 5 -> 1 (since 5 * 3 = 15 = 1)
    \draw[->, trapred, line width=1.4pt] (\aFive:\R-0.15) -- node[midway, below left, font=\tiny\bfseries, trapred, fill=white, inner sep=1.2pt]{$\times 3$} (\aOne:\R-0.15);

    % Highlight 3 and 5 as inverse pair
    \draw[trapred, thick] (\aThree:\R+0.25) circle (4.5pt);
    \draw[trapred, thick] (\aFive:\R+0.25) circle (4.5pt);

    \node[anchor=north, font=\tiny, trapred!90!black, align=center] at (0, -1.85)
      {$3 \times 5 = 15 \equiv 1 \pmod 7$\\$\implies 3^{-1} \equiv 5 \pmod 7$};
  \end{scope}
\end{tikzpicture}
\caption{Modular arithmetic on cyclic dials. (a) Modular addition in $\Z_{12}$:
stepping forward $7$ units from $9$ wraps past the origin and lands on $4$,
discarding full rotations. (b) Modular multiplication and inverses in $\Z_7$:
multiplying by $3$ permutes the nonzero residues in a cyclic orbit
$1 \to 3 \to 2 \to 6 \to 4 \to 5 \to 1$. Because $3 \times 5 = 15 \equiv 1 \pmod 7$,
the residues $3$ and $5$ are multiplicative inverses of each other.}
\label{ch30:fig-clock}
\end{figure}

\begin{definition}[Congruence modulo $n$]
\label{ch30:def-congruence}
For a positive integer $n$, we say $a$ is \emph{congruent to} $b$ \emph{modulo}
$n$, written $a \equiv b \pmod{n}$, if $n \mid (a - b)$; equivalently, $a$ and
$b$ leave the same remainder when divided by $n$.
\end{definition}

So $9 + 7 = 16 \equiv 4 \pmod{12}$. Congruence is an \emph{equivalence relation}
-- reflexive, symmetric, transitive -- and it carves the integers into $n$
buckets called \emph{residue classes}. The class of $a$ is written $[a]_n$ and
collects every integer that lands on the same dial position as $a$. We name the
set of classes $\Z_n = \{[0]_n, [1]_n, \dots, [n-1]_n\}$, and in practice we just
use the representatives $\{0,1,\dots,n-1\}$.

The key property, which justifies all our algorithms, is that congruence
\emph{respects} addition and multiplication.

\begin{theorem}[Arithmetic respects congruence]
\label{ch30:thm-respect}
If $a \equiv a' \pmod n$ and $b \equiv b' \pmod n$, then
\[
  a + b \equiv a' + b' \pmod n
  \qquad\text{and}\qquad
  a\,b \equiv a'\,b' \pmod n.
\]
\end{theorem}

In practice this means \emph{you may reduce mod $n$ at any time}.
To compute $(123 \cdot 456) \bmod 7$ you do not need the product $56088$. Reduce
first: $123 \equiv 4$ and $456 \equiv 1 \pmod 7$, so the answer is $4 \cdot 1 = 4
\bmod 7$. This ``reduce early, reduce often'' habit is what keeps modular
algorithms from ever building giant intermediate numbers.

\begin{trapbox}
Addition and multiplication carry over to the clock perfectly. \emph{Division
does not.} You cannot in general ``cancel'' a common factor: from $6 \equiv 0
\pmod 6$ you cannot conclude $2 \equiv 0 \pmod 6$. Cancelling $a$ from both sides
of a congruence is legal only when $\gcd(a,n) = 1$. Getting this wrong is the
classic modular-arithmetic blunder, and untangling exactly when division
\emph{is} allowed is the subject of the next section.
\end{trapbox}

\subsection{Units and inverses}

On the clock $\Z_n$, the element $1$ acts as a multiplicative identity. An
element $a$ has a \emph{multiplicative inverse} if there is some $a^{-1}$ with
$a \cdot a^{-1} \equiv 1 \pmod n$. When such an inverse exists we call $a$ a
\emph{unit}. Inverses are what let us ``divide'': dividing by $a$ means
multiplying by $a^{-1}$.

\begin{theorem}[Existence of inverses]
\label{ch30:thm-inverse}
The element $a$ has a multiplicative inverse modulo $n$ if and only if
$\gcd(a,n) = 1$. When it exists, the inverse is unique modulo $n$.
\end{theorem}

\begin{intuitionbox}
This follows immediately from B\'ezout's identity. When $\gcd(a,n) = 1$,
the extended Euclidean algorithm returns integers $x, y$ with $ax + ny = 1$.
Reading that equation modulo $n$ drops the $ny$ term (a multiple of $n$) and
leaves $a x \equiv 1 \pmod n$, so $x \bmod n$ \emph{is} the inverse. Conversely, if $\gcd(a,n) = g > 1$, then every multiple of $a$ is a
multiple of $g$, so $ax$ can never be $\equiv 1$; no inverse exists.
\end{intuitionbox}

The set of units modulo $n$ is written $\Z_n^{*} = \{ a : 1 \le a \le n-1,\;
\gcd(a,n) = 1 \}$, and its size has a name: $\phi(n)$, \emph{Euler's totient}.
For a prime $p$, every nonzero element is a unit, so $\phi(p) = p - 1$. For a
product of two distinct primes, $\phi(pq) = (p-1)(q-1)$ -- a fact we will lean on
hard when we build RSA.

\begin{example}[Inverting $5$ modulo $12$]
\label{ch30:ex-inverse}
Since $\gcd(5,12) = 1$, an inverse exists. Run extended Euclid on $(5,12)$ or
$(12,5)$: we find $5 \cdot 5 + 12 \cdot (-2) = 25 - 24 = 1$, so $5^{-1} \equiv 5
\pmod{12}$. Check: $5 \cdot 5 = 25 \equiv 1 \pmod{12}$. (Some elements are their
own inverse; $5$ here is one of them.) By contrast $6$ has \emph{no} inverse mod
$12$ because $\gcd(6,12) = 6 \neq 1$: $6 \cdot k$ runs through $6, 0, 6, 0,
\dots$ and never hits $1$.
\end{example}

\section{Solving Modular Linear Equations}
\label{ch30:linear}

We can now solve the general modular linear congruence: find all $x$ satisfying
\[
  a x \equiv b \pmod n .
\]
This is the modular cousin of the line $ax = b$. The number of solutions is
governed entirely by $\gcd(a,n)$.

\begin{theorem}[Solutions of $ax \equiv b \pmod n$]
\label{ch30:thm-linear}
Let $g = \gcd(a,n)$. The congruence $ax \equiv b \pmod n$ has a solution if and
only if $g \mid b$. When $g \mid b$, there are exactly $g$ solutions modulo $n$,
and they are evenly spaced $n/g$ apart.
\end{theorem}

\begin{intuitionbox}
Why $g \mid b$? As $x$ ranges over all integers, $ax \bmod n$ only ever lands on
multiples of $g$ -- that is exactly what B\'ezout told us, that $ax + ny$ can
only produce multiples of $g$. So if $b$ is not a multiple of $g$, the target is
simply unreachable, like trying to land on an odd number while only ever taking
even-sized steps. When $b$ \emph{is} reachable, the symmetry of the clock means
the solutions repeat in a regular pattern, and there turn out to be exactly $g$
of them.
\end{intuitionbox}

The recipe to find them: use extended Euclid to get $x', y'$ with $ax' + ny' =
g$. Scale by $b/g$ to get one solution $x_0 = x' \cdot (b/g) \bmod n$. The full
solution set is then $x_0 + k\,(n/g)$ for $k = 0,1,\dots,g-1$.

\begin{algorithm}[H]
\caption{\textsc{Modular-Linear-Solver}$(a,b,n)$}
\label{ch30:alg-linear}
\KwIn{Integers $a,b$ and modulus $n>0$.}
\KwOut{All solutions of $ax\equiv b \pmod n$, or report none.}
$(g,\, x',\, y') \leftarrow \textsc{Extended-Euclid}(a \bmod n,\; n)$\;
\If{$g \nmid b$}{\Return ``no solution''\;}
$x_0 \leftarrow x' \cdot (b / g) \bmod n$\;
\For{$k \leftarrow 0$ \KwTo $g-1$}{
  \textbf{output} $\bigl(x_0 + k\cdot(n/g)\bigr) \bmod n$\;
}
\end{algorithm}

\begin{example}[Solving $14x \equiv 30 \pmod{100}$]
\label{ch30:ex-linear}
Here $g = \gcd(14,100) = 2$, and $2 \mid 30$, so solutions exist -- exactly $2$
of them. Extended Euclid gives $14 \cdot (-7) + 100 \cdot 1 = 2$, so $x' = -7$.
Scale by $b/g = 15$: $x_0 = (-7)\cdot 15 \bmod 100 = -105 \bmod 100 = 95$. The
two solutions, spaced $n/g = 50$ apart, are
\[
  x \equiv 95 \quad\text{and}\quad x \equiv 95 + 50 = 145 \equiv 45 \pmod{100}.
\]
Check $x = 45$: $14 \cdot 45 = 630 = 6\cdot 100 + 30 \equiv 30 \pmod{100}$.
And $x = 95$: $14 \cdot 95 = 1330 = 13\cdot 100 + 30 \equiv 30$. Both land.
\end{example}

When $\gcd(a,n) = 1$ there is a single solution, $x \equiv a^{-1} b$, and this
special case -- one equation, one answer -- is the one we use most often.

\section{The Chinese Remainder Theorem}
\label{ch30:crt}

The next result goes back to an ancient puzzle. A 3rd-century
Chinese text asks: \emph{there is a number which leaves remainder $2$ when
divided by $3$, remainder $3$ when divided by $5$, and remainder $2$ when divided
by $7$ -- what is it?} The Chinese Remainder Theorem (CRT) says such a number
always exists, is unique modulo $3\cdot5\cdot7 = 105$, and tells you how to find
it.

\begin{theorem}[Chinese Remainder Theorem]
\label{ch30:thm-crt}
Let $n_1, n_2, \dots, n_k$ be \emph{pairwise coprime} positive integers and let
$N = n_1 n_2 \cdots n_k$. Then for any target remainders $a_1, \dots, a_k$ the
system
\[
  x \equiv a_1 \pmod{n_1},\quad
  x \equiv a_2 \pmod{n_2},\quad \dots,\quad
  x \equiv a_k \pmod{n_k}
\]
has a unique solution $x$ modulo $N$.
\end{theorem}

\begin{intuitionbox}
The simplest way to understand CRT is as a \emph{change of coordinates}. A single
number $x$ in the range $0 \le x < N$ can be described in two completely
equivalent ways: by its ordinary value, or by the tuple of its remainders
$(x \bmod n_1, \dots, x \bmod n_k)$. CRT says this tuple is a perfect ID badge:
no two numbers in $[0,N)$ share a badge, and every possible badge belongs to
exactly one number. So computing with big numbers mod $N$ is the same as
computing with several small numbers in parallel, one per coordinate -- and
addition and multiplication act coordinate by coordinate. It is the integer
version of breaking a vector into its components.
\end{intuitionbox}

Figure~\ref{ch30:fig-crt} draws the bijection for the small case $N = 3\cdot 5 =
15$: each of the $15$ numbers gets a unique two-coordinate badge $(x \bmod 3,\,
x \bmod 5)$, and every cell of the $3 \times 5$ grid is filled exactly once.

\begin{figure}[H]
\centering
\begin{tikzpicture}[scale=0.62]
  % grid columns = x mod 5 (0..4), rows = x mod 3 (0..2)
  \foreach \c in {0,...,4}{ \node[font=\scriptsize] at (\c+0.5,3.4) {\c}; }
  \foreach \r in {0,...,2}{ \node[font=\scriptsize] at (-0.5,2.5-\r) {\r}; }
  \node[font=\scriptsize] at (2.5,4.05) {$x \bmod 5$};
  \node[font=\scriptsize, rotate=90] at (-1.15,1.5) {$x \bmod 3$};
  % highlight x=23 mod 15 = 8 -> (2,3)
  \fill[targetgreen!25] (3,0) rectangle (4,1);
  \draw (0,0) grid (5,3);
  \foreach \x in {0,...,14}{
    \pgfmathtruncatemacro\rr{mod(\x,3)}
    \pgfmathtruncatemacro\cc{mod(\x,5)}
    \ifnum\x=8
      \node[font=\small\bfseries, targetgreen!80!black] at (\cc+0.5, 2.5-\rr) {\x};
    \else
      \node[font=\small] at (\cc+0.5, 2.5-\rr) {\x};
    \fi
  }
\end{tikzpicture}
\caption{CRT as a change of coordinates, $N=15$. Row $=x\bmod 3$, column $=x
\bmod 5$. Every number $0\le x<15$ lands in a distinct cell, and every cell is
occupied: the map $x \mapsto (x\bmod 3,\,x\bmod 5)$ is a bijection. The
highlighted cell $(2,3)$ holds the unique $x=8$.}
\label{ch30:fig-crt}
\end{figure}

\subsection{Constructing the solution}

For each $i$, let
$M_i = N / n_i$ -- the product of \emph{all the other} moduli. The point of
$M_i$ is that it is divisible by every $n_j$ except $n_i$, so it is invisible
(zero) in every coordinate but the $i$-th. Because the moduli are pairwise
coprime, $\gcd(M_i, n_i) = 1$, so $M_i$ has an inverse $M_i^{-1} \pmod{n_i}$.
Define $c_i = M_i \cdot (M_i^{-1} \bmod n_i)$. Then $c_i \equiv 1 \pmod{n_i}$ and
$c_i \equiv 0 \pmod{n_j}$ for $j \neq i$ -- an ``indicator'' for
coordinate $i$. The answer is the weighted sum
\[
  x \equiv \sum_{i=1}^{k} a_i\, c_i \pmod N .
\]
Each term contributes $a_i$ to its own coordinate and $0$ to every other, so the
whole sum has exactly the remainders we ordered.

\begin{algorithm}[H]
\caption{\textsc{CRT}$(a_1,\dots,a_k;\, n_1,\dots,n_k)$}
\label{ch30:alg-crt}
\KwIn{Remainders $a_i$ and pairwise-coprime moduli $n_i$.}
\KwOut{The unique $x \in [0,N)$ solving all congruences, $N = \prod n_i$.}
$N \leftarrow n_1 n_2 \cdots n_k$\;
$x \leftarrow 0$\;
\For{$i \leftarrow 1$ \KwTo $k$}{
  $M_i \leftarrow N / n_i$\;
  $(g, s, t) \leftarrow \textsc{Extended-Euclid}(M_i \bmod n_i,\; n_i)$\;
  \tcp{$s$ is $M_i^{-1} \bmod n_i$ since $g=1$}
  $x \leftarrow \bigl(x + a_i \cdot M_i \cdot s\bigr) \bmod N$\;
}
\Return $x$\;
\end{algorithm}

\begin{example}[The ancient puzzle]
\label{ch30:ex-crt}
Solve $x \equiv 2 \pmod 3$, $x \equiv 3 \pmod 5$, $x \equiv 2 \pmod 7$. Here
$N = 105$.
\begin{center}
\begin{tabular}{c c c c c c}
\toprule
$i$ & $n_i$ & $a_i$ & $M_i = N/n_i$ & $M_i \bmod n_i$ & $M_i^{-1} \bmod n_i$ \\
\midrule
$1$ & $3$ & $2$ & $35$ & $2$ & $2$ \\
$2$ & $5$ & $3$ & $21$ & $1$ & $1$ \\
$3$ & $7$ & $2$ & $15$ & $1$ & $1$ \\
\bottomrule
\end{tabular}
\end{center}
Now sum the weighted terms:
\[
  x \equiv 2\cdot 35\cdot 2 \;+\; 3\cdot 21\cdot 1 \;+\; 2\cdot 15\cdot 1
    = 140 + 63 + 30 = 233 \equiv 233 - 2\cdot 105 = 23 \pmod{105}.
\]
So $x = 23$. Verify: $23 = 7\cdot3 + 2 \equiv 2 \pmod 3$; $23 = 4\cdot5 + 3
\equiv 3 \pmod 5$; $23 = 3\cdot7 + 2 \equiv 2 \pmod 7$. All three hold.
\end{example}

\begin{aestheticbox}
The Chinese Remainder Theorem also speeds up RSA decryption by a factor of roughly
four: instead of one exponentiation modulo the big number $N = pq$, you do two
exponentiations modulo the half-size primes $p$ and $q$ and recombine with CRT.
Since modular exponentiation cost grows faster than linearly in the bit length,
two half-size jobs beat one full-size job. CRT also underlies fast big-integer
arithmetic, error-correcting codes, and hash-table tricks. All of these use the same idea of
computing in parallel coordinates.
\end{aestheticbox}

\section{Modular Exponentiation}
\label{ch30:modexp}

RSA needs to compute things like $a^{b} \bmod n$ where $a$, $b$, and $n$ are all
hundreds of digits long. Two naive approaches both fail. Multiplying
$a$ by itself $b$ times takes $b$ multiplications -- but $b$ might be $2^{1024}$,
more multiplications than there are atoms in the universe. And computing the full
$a^b$ first and reducing at the end is worse still: $a^b$ has astronomically many
digits and would never fit in memory. We need an algorithm that is fast in the
number of multiplications \emph{and} keeps every intermediate value small.

Both problems are solved by one idea: \emph{repeated squaring}, also called
\emph{binary exponentiation}.

\begin{intuitionbox}
To compute $a^{13}$, do not multiply thirteen copies of $a$. Instead notice that
$13$ in binary is $1101$, i.e.\ $13 = 8 + 4 + 1$, so
$a^{13} = a^{8}\cdot a^{4}\cdot a^{1}$. And the powers $a^1, a^2, a^4, a^8$ form a
chain where each is the \emph{square} of the previous: $a^2 = (a^1)^2$, $a^4 =
(a^2)^2$, $a^8 = (a^4)^2$. So you compute the chain by repeated squaring -- just
a handful of multiplications -- and multiply in the ones that correspond to a
$1$-bit of the exponent. The bits of the exponent are a recipe, read one at a
time. An exponent with $\beta$ bits needs only about $\beta$ squarings, not
$2^{\beta}$ multiplications.
\end{intuitionbox}

The clean way to implement this scans the bits of $b$ from the \emph{most}
significant down to the least. Maintain a running result $r$, initially $1$. For
each bit, square $r$ (this ``shifts'' the partial exponent left by one binary
place); if the current bit is $1$, also multiply by $a$. Reduce mod $n$ after
every operation, so nothing ever exceeds $n^2$ in size.

\begin{algorithm}[H]
\caption{\textsc{Modular-Power}$(a,b,n)$ -- computes $a^b \bmod n$}
\label{ch30:alg-modexp}
\KwIn{Base $a$, exponent $b \ge 0$ with bits $b_{t-1} b_{t-2}\cdots b_1 b_0$, modulus $n>1$.}
\KwOut{$a^b \bmod n$.}
$r \leftarrow 1$\;
\For{$i \leftarrow t-1$ \KwDownTo $0$}{
  $r \leftarrow (r \cdot r) \bmod n$\;
  \If{$b_i = 1$}{
    $r \leftarrow (r \cdot a) \bmod n$\;
  }
}
\Return $r$\;
\end{algorithm}

\begin{example}[Computing $7^{13} \bmod 11$]
\label{ch30:ex-modexp}
The exponent $13 = 1101_2$, so we scan bits $1,1,0,1$ from the top. Start
$r = 1$.
\begin{center}
\begin{tabular}{c c l l}
\toprule
bit $b_i$ & operation & $r$ after square & $r$ after multiply \\
\midrule
$1$ & sq, mul & $1^2 = 1$              & $1\cdot 7 = 7$ \\
$1$ & sq, mul & $7^2 = 49 \equiv 5$    & $5\cdot 7 = 35 \equiv 2$ \\
$0$ & sq      & $2^2 = 4$              & (skip) $\;4$ \\
$1$ & sq, mul & $4^2 = 16 \equiv 5$    & $5\cdot 7 = 35 \equiv 2$ \\
\bottomrule
\end{tabular}
\end{center}
Figure~\ref{ch30:fig-modexp} shows how the bits of the exponent drive the
square/multiply chain: every bit triggers a square, and a $1$-bit triggers an
extra multiply by the base.

\begin{figure}[H]
\centering
\begin{tikzpicture}[
  node distance=4mm,
  bit/.style={draw, minimum size=6mm, font=\small},
  op/.style={font=\scriptsize, align=center},
  ar/.style={-{Stealth}, thick}]
  % exponent bits of 13 = 1101, MSB first
  \node[bit, fill=quotientorange!20] (b3) {1};
  \node[bit, right=of b3, fill=quotientorange!20] (b2) {1};
  \node[bit, right=of b2] (b1) {0};
  \node[bit, right=of b1, fill=quotientorange!20] (b0) {1};
  \node[left=5mm of b3, font=\scriptsize] {$13=$};
  % operations under each bit
  \node[op, below=6mm of b3] (o3) {square\\ $\times a$};
  \node[op, below=6mm of b2] (o2) {square\\ $\times a$};
  \node[op, below=6mm of b1] (o1) {square};
  \node[op, below=6mm of b0] (o0) {square\\ $\times a$};
  \draw[ar] (b3) -- (o3); \draw[ar] (b2) -- (o2);
  \draw[ar] (b1) -- (o1); \draw[ar] (b0) -- (o0);
  % running r values
  \node[op, below=4mm of o3, targetgreen] {$r{=}7$};
  \node[op, below=4mm of o2, targetgreen] {$r{=}2$};
  \node[op, below=4mm of o1, targetgreen] {$r{=}4$};
  \node[op, below=4mm of o0, targetgreen] {$r{=}2$};
\end{tikzpicture}
\caption{Bits of the exponent $13=1101_2$ as a recipe for $7^{13}\bmod 11$, read
left to right. Each bit squares the running value $r$; each shaded $1$-bit also
multiplies by the base $7$. The final $r=2$ is the answer.}
\label{ch30:fig-modexp}
\end{figure}

All reductions are mod $11$. The result is $7^{13} \equiv 2 \pmod{11}$. (Sanity
check by Fermat, coming next: $7^{10}\equiv 1$, so $7^{13}\equiv 7^3 = 343 = 31
\cdot 11 + 2 \equiv 2$. Agreed.) Notice every intermediate value stayed below
$121 = 11^2$; we never built a giant number.
\end{example}

\begin{complexitybox}
\textbf{Modular exponentiation.} For a $\beta$-bit exponent $b$ and modulus $n$:
\begin{itemize}
  \item \emph{Multiplications:} between $\beta$ and $2\beta$ modular
        multiplications (one square per bit, plus one multiply per $1$-bit).
  \item \emph{Intermediate size:} every value stays in $[0,n)$, so each multiply
        is of $\beta$-bit numbers, costing $O(\beta^2)$ bit operations naively.
  \item \emph{Total bit cost:} $O(\beta^3)$ with schoolbook multiplication.
\end{itemize}
Contrast the naive $b$ multiplications, which is $\Theta(2^{\beta})$; for
exponents of cryptographic size the naive method would never finish.
\end{complexitybox}

\begin{walkbox}
Take \textsc{Modular-Power} on a walk. Convince yourself of its loop invariant:
just before processing bit $b_i$, the value of $r$ equals $a^{(\text{bits seen so
far})} \bmod n$, where ``bits seen so far'' is the number formed by the
high-order bits $b_{t-1}\cdots b_{i+1}$ already consumed. Squaring doubles that
prefix value (appends a $0$ bit); a conditional multiply by $a$ turns the new
last bit into a $1$. When the loop ends, every bit has been folded in, and $r =
a^b \bmod n$. Why is squaring exactly ``append a zero bit'' to the exponent?
\end{walkbox}

\subsection{Fermat's and Euler's theorems}

Two classical facts about exponents underpin both the primality test and RSA.

\begin{theorem}[Fermat's little theorem]
\label{ch30:thm-fermat}
If $p$ is prime and $\gcd(a,p) = 1$, then $a^{p-1} \equiv 1 \pmod p$.
\end{theorem}

\begin{theorem}[Euler's theorem]
\label{ch30:thm-euler}
If $\gcd(a,n) = 1$, then $a^{\phi(n)} \equiv 1 \pmod n$, where $\phi$ is Euler's
totient.
\end{theorem}

Euler's theorem generalizes Fermat's (take $n = p$ prime, so $\phi(p) = p-1$).
The intuition for both: the units mod $n$ form a finite group under
multiplication, and in any finite group, raising any element to the size of the
group lands you back at the identity. We will use Euler's theorem directly to
prove RSA decryption inverts encryption.

\section{The RSA Public-Key Cryptosystem}
\label{ch30:rsa}

RSA, named for Rivest, Shamir,
and Adleman, lets two strangers who have never met exchange a secret over an open
channel. The key idea is that the encryption key can be published,
while only the holder of a separate private key can
decrypt. This is \emph{public-key} cryptography, and before 1977 most people
thought it was impossible.

\begin{intuitionbox}
Think of a padlock that anyone can snap shut but only you have the key to open.
You manufacture thousands of these open padlocks and scatter them everywhere -- a
mailbox at every street corner. Anyone can grab one, lock a box with it, and send
you the box; nobody but you can reopen it, because only you hold the key. The
public key is the open padlock; the private key is the key in your pocket. The
mathematical question is: how do you build a ``lock'' that is easy to snap shut
(encrypt) but, without the secret, ruinously hard to open (decrypt)? RSA's answer
is the gap between multiplying primes (easy) and factoring their product (hard).
\end{intuitionbox}

\subsection{Setting up the keys}

\begin{enumerate}
  \item Pick two large random primes $p$ and $q$ (in practice, each several
        hundred digits). Compute $n = pq$.
  \item Compute $\phi(n) = (p-1)(q-1)$.
  \item Choose a public exponent $e$ with $1 < e < \phi(n)$ and $\gcd(e,
        \phi(n)) = 1$. A common choice is $e = 65537$.
  \item Compute the private exponent $d = e^{-1} \bmod \phi(n)$ using extended
        Euclid -- this is the one step that needs the secret factorization.
\end{enumerate}

The \emph{public key} is the pair $(e, n)$. The \emph{private key} is $d$ (and
implicitly the factors $p, q$). The primes $p$ and $q$, and $\phi(n)$, are
destroyed or hidden after setup.

\subsection{Encrypting and decrypting}

A message is encoded as a number $m$ with $0 \le m < n$. Then:
\[
  \text{encrypt:}\quad c = m^{e} \bmod n,
  \qquad\qquad
  \text{decrypt:}\quad m = c^{d} \bmod n.
\]
Figure~\ref{ch30:fig-rsa} traces a single message across the channel.

\begin{figure}[H]
\centering
\begin{tikzpicture}[
  node distance=7mm and 12mm,
  box/.style={draw, rounded corners, align=center, font=\small, minimum height=9mm},
  pub/.style={box, fill=sourceblue!10},
  prv/.style={box, fill=targetgreen!12},
  flow/.style={-{Stealth}, thick}]
  \node[box] (m) {message\\$m$};
  \node[pub, right=of m] (enc) {encrypt\\$c=m^{e}\bmod n$};
  \node[box, right=of enc] (c) {ciphertext\\$c$};
  \node[prv, right=of c] (dec) {decrypt\\$m=c^{d}\bmod n$};
  \node[box, right=of dec] (m2) {message\\$m$};
  \draw[flow] (m) -- (enc);
  \draw[flow] (enc) -- (c);
  \draw[flow] (c) -- (dec);
  \draw[flow] (dec) -- (m2);
  \node[sourceblue, font=\scriptsize, below=2mm of enc] {public key $(e,n)$};
  \node[targetgreen, font=\scriptsize, below=2mm of dec] {private key $d$};
  \node[font=\scriptsize, above=1mm of c] {sent over open channel};
\end{tikzpicture}
\caption{The RSA round trip. Anyone can encrypt with the published $(e,n)$;
only the holder of the private $d$ can turn the ciphertext back into $m$. Both
arrows are a single call to \textsc{Modular-Power}.}
\label{ch30:fig-rsa}
\end{figure}

Both are just \textsc{Modular-Power}. Encryption uses only the public $(e,n)$;
decryption needs the private $d$. The whole security rests on this asymmetry:
to recover $d$ from the public $(e,n)$ you would need $\phi(n) = (p-1)(q-1)$,
which requires knowing $p$ and $q$ -- that is, factoring $n$. For large $n$,
nobody knows how to do that in any reasonable time.

\begin{theorem}[RSA correctness]
\label{ch30:thm-rsa}
With keys set up as above, decryption inverts encryption: $(m^e)^d \equiv m
\pmod n$ for every $m$ with $0 \le m < n$.
\end{theorem}

\begin{formalbox}
\textit{Proof sketch.} By construction $ed \equiv 1 \pmod{\phi(n)}$, so $ed = 1 +
k\,\phi(n)$ for some integer $k$. Hence
\[
  m^{ed} = m^{1 + k\phi(n)} = m \cdot \bigl(m^{\phi(n)}\bigr)^{k}.
\]
When $\gcd(m,n) = 1$, Euler's theorem (Theorem~\ref{ch30:thm-euler}) gives
$m^{\phi(n)} \equiv 1 \pmod n$, so the whole second factor is $1$ and $m^{ed}
\equiv m \pmod n$. The remaining cases, where $m$ shares a factor with $n$, are
handled by applying Fermat's little theorem modulo $p$ and modulo $q$ separately
and stitching the two results back together with the Chinese Remainder Theorem;
the congruence $m^{ed} \equiv m$ holds modulo both primes, hence modulo their
product $n$. \hfill$\square$
\end{formalbox}

\begin{example}[A toy RSA, end to end]
\label{ch30:ex-rsa}
Use tiny primes so we can do everything by hand. Let $p = 5$, $q = 11$, so $n =
55$ and $\phi(n) = 4 \cdot 10 = 40$. Pick $e = 3$; check $\gcd(3,40) = 1$, good.
Find $d = 3^{-1} \bmod 40$: extended Euclid gives $3 \cdot 27 = 81 = 2\cdot 40 +
1$, so $d = 27$.

Encrypt the message $m = 7$:
\[
  c = 7^{3} \bmod 55 = 343 \bmod 55 = 343 - 6\cdot 55 = 13.
\]
Decrypt $c = 13$ with $d = 27$:
\[
  m = 13^{27} \bmod 55.
\]
Compute by repeated squaring mod $55$: $13^2 = 169 \equiv 4$, $13^4 \equiv 16$,
$13^8 \equiv 256 \equiv 36$, $13^{16} \equiv 36^2 = 1296 \equiv 31$. Since $27 =
16 + 8 + 2 + 1$,
\[
  13^{27} \equiv 31 \cdot 36 \cdot 4 \cdot 13 \pmod{55}.
\]
Multiply step by step mod $55$: $31\cdot 36 = 1116 \equiv 1116 - 20\cdot 55 = 16$;
$16 \cdot 4 = 64 \equiv 9$; $9 \cdot 13 = 117 \equiv 117 - 2\cdot 55 = 7$. We
recover $m = 7$. The round trip worked.
\end{example}

\begin{aestheticbox}
This toy is trivial to break -- anyone can factor $55 = 5\cdot 11$. Real RSA
uses $n$ with around $2048$ to $4096$ bits, where the best known factoring
algorithms would run for longer than the age of the universe. Every time your
browser shows a padlock icon, a key exchange related to ideas in this chapter
just happened. The entire scheme is built from the
elementary tools we developed by hand: GCD, modular inverses, modular
exponentiation, and Euler's theorem.
\end{aestheticbox}

\begin{trapbox}
Textbook RSA, exactly as described, is \emph{not} safe to deploy. It is
deterministic (the same message always encrypts to the same ciphertext, leaking
information), it is malleable, and small messages with small $e$ can be attacked
directly. Real systems add randomized padding (OAEP) and use RSA only to wrap a
symmetric key. Learn the math here; never ship the textbook version.
\end{trapbox}

\section{Primality Testing}
\label{ch30:primality}

RSA needs large primes, hundreds of digits each. So we need to answer, quickly:
\emph{is this 300-digit number prime?} Trial division -- checking every divisor
up to $\sqrt n$ -- would require about $10^{150}$ operations for a $300$-digit
number, which is hopeless. We need something exponentially faster, and
Fermat's little theorem, used \emph{in reverse}, provides it.

\subsection{The Fermat test and its flaw}

Fermat's little theorem says: if $p$ is prime and $\gcd(a,p)=1$, then $a^{p-1}
\equiv 1 \pmod p$. Turn it into a test. To check whether $n$ is prime, pick a
random ``witness'' $a$ and compute $a^{n-1} \bmod n$ by repeated squaring. If
the result is \emph{not} $1$, then $n$ is \emph{definitely composite}, since
Fermat's little theorem rules this out for every prime. If the result \emph{is}
$1$, then $n$ is \emph{probably prime}, and we gain confidence by trying more
witnesses.

\begin{trapbox}
The word ``probably'' matters. Passing for a single base does not prove
primality: some composites pass for specific bases (making such a base a
\emph{Fermat liar}), and \emph{Carmichael numbers} (the smallest is
$561 = 3\cdot 11\cdot 17$) satisfy $a^{n-1}\equiv 1 \pmod n$ for \emph{every}
$a$ coprime to $n$, fooling the plain Fermat test on every coprime base.
Carmichael numbers are why we need a stronger test.
\end{trapbox}

\subsection{Miller--Rabin: a better witness}

Miller--Rabin patches the hole with one extra observation. If $n$ is an odd prime,
then in $\Z_n$ the only square roots of $1$ are $+1$ and $-1$. (In a field, $x^2 -
1 = (x-1)(x+1) = 0$ forces $x = \pm 1$; composites can have \emph{nontrivial} square
roots of $1$, which immediately prove compositeness.) Miller--Rabin watches the squaring
chain for an illegal square root of $1$.

Write $n - 1 = 2^{s}\cdot u$ with $u$ odd -- just pull out
all the factors of two. Now look at the sequence
\[
  a^{u},\;\; a^{2u},\;\; a^{4u},\;\; \dots,\;\; a^{2^{s} u} = a^{n-1} \pmod n,
\]
each term the square of the one before it. If $n$ is prime, this sequence must
either start at $1$, or hit $-1$ at some point before the end (and stay $1$
thereafter). If instead the sequence reaches $1$ \emph{without} the step just
before it being $-1$, then we have found a nontrivial square root of $1$,
proving that $n$ is composite. We call such a base $a$ a \emph{witness to
compositeness}.

\begin{algorithm}[H]
\caption{\textsc{Miller--Rabin-Witness}$(a,n)$ -- does $a$ prove $n$ composite?}
\label{ch30:alg-mr}
\KwIn{Base $a$ with $2 \le a \le n-2$; odd $n > 2$.}
\KwOut{\textsc{true} if $a$ proves $n$ composite; \textsc{false} if $n$ passes for $a$.}
write $n - 1 = 2^{s}\cdot u$ with $u$ odd\;
$x \leftarrow \textsc{Modular-Power}(a, u, n)$\;
\If{$x = 1$ \textbf{or} $x = n-1$}{\Return \textsc{false}\;}
\For{$j \leftarrow 1$ \KwTo $s-1$}{
  $x \leftarrow (x \cdot x) \bmod n$\;
  \If{$x = n-1$}{\Return \textsc{false}\;}
}
\Return \textsc{true}\;
\end{algorithm}

To test $n$, run this witness check for several random bases $a$. If any base
returns \textsc{true}, $n$ is composite, with certainty. If all of them return
\textsc{false}, declare $n$ ``probably prime.''

\begin{theorem}[Miller--Rabin error bound]
\label{ch30:thm-mr}
If $n$ is an odd composite, then at least $3/4$ of the bases $a \in \{2, \dots,
n-2\}$ are witnesses to its compositeness. Hence a single random base fails to
detect a composite with probability at most $1/4$, and $k$ independent random
bases reduce the failure probability to at most $4^{-k}$.
\end{theorem}

\begin{intuitionbox}
The $4^{-k}$ decay is what makes this practical. With $k = 50$ random bases the
chance of wrongly calling a composite ``prime'' is below $4^{-50} \approx
10^{-30}$ -- vastly smaller than the chance of a cosmic ray flipping a bit in
your computer mid-calculation. For all engineering purposes that is certainty.
And unlike the Fermat test, no composite -- not even a Carmichael number -- can
fool a $3/4$ fraction of bases, since Theorem~\ref{ch30:thm-mr} holds for every
odd composite. We accept a tiny error probability in exchange for a huge speedup.
\end{intuitionbox}

\begin{example}[Catching $n = 561$, the smallest Carmichael number]
\label{ch30:ex-mr}
The Fermat test fails on $561$, so let us watch Miller--Rabin succeed. Here
$n - 1 = 560 = 2^{4}\cdot 35$, so $s = 4$ and $u = 35$. Take base $a = 2$.
First $2^{35} \bmod 561$. By repeated squaring one finds $2^{35} \equiv 263
\pmod{561}$. That is neither $1$ nor $560$, so we begin squaring:
\[
  263^2 \equiv 166,\quad 166^2 \equiv 67,\quad 67^2 \equiv 1 \pmod{561}.
\]
The sequence hit $1$, but the term just before it was $67$, \emph{not} $560 \equiv
-1$. We have caught a nontrivial square root of $1$: $67^2 \equiv 1$ yet $67 \neq
\pm 1$. So $a = 2$ is a witness, and $561$ is declared composite -- correctly,
since $561 = 3 \cdot 11 \cdot 17$. The extra square-root check is exactly what the
plain Fermat test was missing.
\end{example}

Figure~\ref{ch30:fig-mr} contrasts the two endings of the squaring chain. A
prime is allowed to reach $1$ only by passing through $-1$ first; the composite
$561$ jumps directly to $1$ from $67$, a nontrivial square root of $1$ that proves $561$ is composite.

\begin{figure}[H]
\centering
\begin{tikzpicture}[
  node distance=10mm,
  val/.style={draw, circle, minimum size=8.5mm, inner sep=0pt, font=\scriptsize},
  good/.style={val, fill=targetgreen!15},
  bad/.style={val, fill=trapred!18},
  sq/.style={-{Stealth}, thick}]
  % top row: a legal prime ending ... -> -1 -> 1
  \node[font=\scriptsize] at (-2.2,0) {prime, OK:};
  \node[val] (p0) at (0,0) {$x$};
  \node[val] (p1) at (1.6,0) {$\cdots$};
  \node[good] (p2) at (3.2,0) {$-1$};
  \node[good] (p3) at (4.8,0) {$1$};
  \draw[sq] (p0) -- (p1); \draw[sq] (p1) -- (p2); \draw[sq] (p2) -- (p3);
  \node[font=\tiny, targetgreen] at (4.0,0.55) {$(-1)^2=1$};
  % bottom row: composite 561 ... 67 -> 1 (illegal)
  \node[font=\scriptsize] at (-2.2,-1.5) {$561$, caught:};
  \node[val] (q0) at (0,-1.5) {$263$};
  \node[val] (q1) at (1.6,-1.5) {$166$};
  \node[bad] (q2) at (3.2,-1.5) {$67$};
  \node[bad] (q3) at (4.8,-1.5) {$1$};
  \draw[sq] (q0) -- (q1); \draw[sq] (q1) -- (q2); \draw[sq] (q2) -- (q3);
  \node[font=\tiny, trapred] at (4.0,-0.95) {$67^2=1$};
  \node[font=\tiny, trapred] at (4.0,-2.18) {$67\neq\pm1$};
\end{tikzpicture}
\caption{Miller--Rabin's square-root check. Each arrow squares (mod $n$). A
prime can only arrive at $1$ from $-1$ (top). The composite $561$ reaches $1$
directly from $67$, a nontrivial square root of $1$ that proves $561$ is composite (bottom).}
\label{ch30:fig-mr}
\end{figure}

\begin{complexitybox}
\textbf{Miller--Rabin primality test.} For an $n$ with $\beta = \log_2 n$ bits,
using $k$ random bases:
\begin{itemize}
  \item Each base runs one modular exponentiation plus up to $s \le \beta$
        squarings: $O(\beta)$ modular multiplications, i.e.\ $O(\beta^3)$ bit
        operations naively.
  \item Total: $O(k\,\beta^3)$ bit operations.
  \item \emph{Error:} at most $4^{-k}$ of declaring a composite prime; never
        errs on a true prime.
\end{itemize}
This is fast enough to generate $2048$-bit primes in milliseconds: pick a random
odd number, run Miller--Rabin, repeat until one passes. The prime number theorem
guarantees primes are dense enough that you do not wait long.
\end{complexitybox}

\begin{historybox}
This material maps to \textbf{CLRS Chapter~31}, ``Number-Theoretic Algorithms.'' Euclid's algorithm appears in Book~VII of the \emph{Elements} (c.~300~BCE), making it the oldest nontrivial algorithm still in daily use. The Chinese Remainder Theorem traces to Sunzi's \emph{Mathematical Classic} (3rd--5th century CE). Fermat stated his little theorem in 1640; Euler proved and generalized it in 1736 with the totient function. Miller's 1976 deterministic test (conditional on the Extended Riemann Hypothesis) was made unconditional and randomized by Rabin in 1980. RSA was published by Rivest, Shamir, and Adleman in 1978; in 1997, British intelligence agency GCHQ revealed that Clifford Cocks, James Ellis, and Malcolm Williamson had discovered equivalent concepts earlier under classified cover. In 2002, Agrawal, Kayal, and Saxena discovered AKS, the first deterministic polynomial-time primality test, though randomized Miller--Rabin remains faster in practice. CLRS develops all of this with more group theory than we have; read it next for the rigorous treatment of $\Z_n^{*}$ as a group and the full proof of the Miller--Rabin error bound.
\end{historybox}

\begin{labbox}
\textbf{Trace it yourself.} Pick up paper and do these by hand before moving on.
\begin{enumerate}
  \item Run \textsc{Euclid}$(89, 55)$ -- consecutive Fibonacci numbers -- and
        count the divisions. You should see the maximum number of steps for
        inputs this size, every quotient equal to $1$ until the very last.
  \item Run \textsc{Extended-Euclid}$(17, 23)$ and confirm the coefficients give
        $17x + 23y = 1$. Then read off $17^{-1} \bmod 23$ from your $x$.
  \item Compute $3^{20} \bmod 100$ with \textsc{Modular-Power}, scanning the
        bits of $20 = 10100_2$. Keep every intermediate value below $100^2$.
  \item Solve the system $x \equiv 1 \pmod 4$, $x \equiv 2 \pmod 9$, $x \equiv 3
        \pmod{25}$ with CRT. (The moduli are pairwise coprime; $N = 900$.)
  \item Run \textsc{Miller--Rabin-Witness} on $n = 221 = 13 \cdot 17$ with base
        $a = 5$ and watch it return \textsc{true}.
\end{enumerate}
\end{labbox}

\section{Putting It Together}
\label{ch30:together}

Every tool in this chapter grew out of a single recurrence from 300~BCE. At the
bottom sits $\gcd(a,b) = \gcd(b, a \bmod b)$, from which comes Euclid's
algorithm. Tracking the quotients along the way yields B\'ezout's coefficients
and modular inverses; inverses solve modular linear equations and drive the
Chinese Remainder Theorem. Alongside them, repeated squaring gives fast modular
exponentiation, Euler's theorem gives the cycle length of modular powers, and
combining modular inverses with Euler's theorem produces RSA -- while adding the
square-root-of-one check gives Miller--Rabin to find the primes RSA needs.

\begin{focusbox}
The algorithm used most in this chapter is \emph{the extended
Euclidean algorithm}. Modular inverses, linear congruences, CRT reconstruction,
and RSA key generation all reduce directly to computing B\'ezout coefficients
with that one routine. Master \textsc{Extended-Euclid}
and \textsc{Modular-Power}, understand why each is fast, and you understand the
computational core of classical cryptography.
\end{focusbox}

The modular arithmetic and prime-modulus hash constructions developed here do not end with cryptosystems. In Chapter~31, we turn from numbers to text: string matching, pattern search, and the Rabin--Karp algorithm, which uses modular hashing to find patterns in long texts in linear expected time.

\section*{Exercises}
\addcontentsline{toc}{section}{Exercises}

\begin{exercise}[Warmup: trace Euclid]
\label{ch30:ex-trace}
Compute $\gcd(1071, 462)$ by hand using \textsc{Euclid}, writing out the
quotient and remainder at each step. How many divisions did it take? Then state
$\gcd$ as a product of shared prime factors to double-check.
\end{exercise}

\begin{exercise}[B\'ezout by hand]
\label{ch30:ex-bezout}
Run \textsc{Extended-Euclid}$(240, 46)$. Report the triple $(d, x, y)$ and verify
$240x + 46y = d$. Is the solution $(x,y)$ unique? If not, describe the full
family of integer solutions.
\end{exercise}

\begin{exercise}[Inverses exist or do not]
\label{ch30:ex-inv}
For each pair, decide whether $a$ has an inverse modulo $n$, and compute it when
it exists: (a) $a=7, n=26$; (b) $a=8, n=12$; (c) $a=10, n=21$. Explain each
answer in terms of $\gcd(a,n)$.
\end{exercise}

\begin{exercise}[Solving congruences]
\label{ch30:ex-cong}
Find all solutions modulo $n$, or explain why none exist:
(a) $6x \equiv 9 \pmod{15}$; \;
(b) $6x \equiv 8 \pmod{15}$; \;
(c) $35x \equiv 14 \pmod{84}$.
For each, first compute $g = \gcd(a,n)$ and use the divisibility test before
hunting for solutions.
\end{exercise}

\begin{exercise}[Chinese Remainder reconstruction]
\label{ch30:ex-crt2}
A basket of eggs breaks into groups with $1$ left over when counted by $3$, $2$
left over by $5$, and $6$ left over by $7$. What is the smallest possible number
of eggs? Use the \textsc{CRT} construction and show your $M_i$ and inverses.
\end{exercise}

\begin{exercise}[Fast exponentiation]
\label{ch30:ex-fastexp}
Compute $5^{117} \bmod 19$ using \textsc{Modular-Power}. Write $117$ in binary,
list the squaring/multiplying steps, and keep every intermediate value reduced
mod $19$. Then verify your answer using Fermat's little theorem to shrink the
exponent first.
\end{exercise}

\begin{exercise}[RSA round trip]
\label{ch30:ex-rsa2}
Let $p = 7$, $q = 13$, so $n = 91$ and $\phi(n) = 72$. Choose $e = 5$.
(a)~Confirm $\gcd(5,72)=1$ and compute the private exponent $d = 5^{-1} \bmod 72$.
(b)~Encrypt $m = 10$. (c)~Decrypt your ciphertext and confirm you recover $10$.
\end{exercise}

\begin{exercise}[Why the square-root test matters]
\label{ch30:ex-mr2}
The number $341 = 11 \cdot 31$ passes the Fermat test for base $a = 2$ (indeed
$2^{340} \equiv 1 \pmod{341}$), so the Fermat test calls it ``probably prime.''
Run \textsc{Miller--Rabin-Witness}$(2, 341)$ instead. Does base $2$ expose $341$
as composite, or does $341$ slip past base $2$ as well? (You will need
$n-1 = 340 = 2^2 \cdot 85$.)
\end{exercise}

\begin{exercise}[Worst case of Euclid]
\label{ch30:ex-fib}
Prove that, for every $k \ge 2$, \textsc{Euclid}$(F_{k+1}, F_k)$ performs exactly $k-1$ division steps,
where $F_i$ is the $i$-th Fibonacci number. (Hint: show by induction that each
reduction turns the pair $(F_{i+1}, F_i)$ into $(F_i, F_{i-1})$ because the
quotient is always $1$.)
\end{exercise}

\begin{exercise}[Challenge: CRT speedup for RSA]
\label{ch30:ex-crtrsa}
Explain, with a rough operation count, why decrypting $c^{d} \bmod n$ by
computing $c^{d \bmod (p-1)} \bmod p$ and $c^{d \bmod (q-1)} \bmod q$ separately
and recombining with CRT is about four times faster than the direct computation
modulo $n = pq$. (Hint: modular multiplication of $\beta$-bit numbers costs
$\Theta(\beta^2)$ naively, and $p$, $q$ each have about half the bits of $n$.)
\end{exercise}

\begin{exercise}[Challenge: a flawed inverse routine]
\label{ch30:ex-bug}
A student computes modular inverses by the formula $a^{-1} \equiv a^{n-2} \pmod
n$, citing Fermat. (a)~For which $n$ is this formula valid? (b)~Give a specific
$a, n$ where it produces a wrong answer, and explain which hypothesis of Fermat's
little theorem fails. (c)~State the correct generalization using Euler's theorem,
and say why \textsc{Extended-Euclid} is preferable in practice anyway.
\end{exercise}
